\scanpage{108}

$K$ extension tot. réelle \margin{[i.e. si $K$ défini par $H\subset\pi_1(Q)$, $H$ ne contient aucun conjugué de $\tau$, $\tau$ la conjugaison ordinaire)}

A) $L$ extension quadratique, $L=K(\unsure{x|z})$, $x^2+ax+b=0$ ($a,b\in K$).
\inserted{$\sigma$ l'élément non trivial de $\operatorname{Gal}(L/K)$, donc $\sigma x=\bar x$ pour $\forall$ $x$}
\margin{$a=x+{}^\sigma x$, $b=x\,{}^\sigma x$}
Conditions équivalentes:
\diagram{\begin{tikzcd}[row sep=large]
\text{$L$ tot. imaginaire}\arrow[d,phantom,"\unsure{\Updownarrow|\Downarrow}"]\\
\begin{gathered}\text{Pour $\forall$ plongt de $K$ dans $R$, }\Delta=\unsure{a^2}-4b\\=(x-{}^\sigma x)^2\text{ est }\unsure{<}0,\text{ i.e. $\Delta$ est}\\\text{tot. \deleted{positif} négatif}\end{gathered}\\
\deleted{\begin{gathered}\text{$L$ est stable par }\lambda\mapsto\bar\lambda,\text{ et}\\(x-\bar x)^2\text{ est tot. positif}\end{gathered}}\arrow[u,Leftrightarrow]
\end{tikzcd}}
\deleted{Alors $\sigma$ est la conj. \unsure{complexe}.}

B) Soit $x\in\bar Q\subset C$. Conditions équivalentes
\diagram{\begin{tikzcd}[row sep=large]
\begin{gathered}1)\quad\left\{\begin{aligned}&x+\bar x\ [\text{et }x\bar x\text{ sont}]\ \text{tot. réels}\\&(x-\bar x)^2\text{ est tot. \deleted{\unsure{positif} \unsure{$\leq0$}} négatif}\end{aligned}\right.\end{gathered}\arrow[d,Leftrightarrow]\\
\begin{gathered}2)\quad\deleted{x\text{ est }\unsure{\text{un}}}\ \exists\text{ un $K$ tot. réel et une}\\\text{extension $L$ tot. imaginaire de $K$, dont}\\\text{$x$ soit un générateur.}\end{gathered}\arrow[d,Leftrightarrow]\\
2')\quad\deleted{\unsure{x}}\ x\in Q'(i)-Q',\quad\text{i.e.}
\end{tikzcd}}
\inserted{conséquence des autres, car $x\bar x=\frac14[(x+\bar x)^2-(x-\bar x)^2]$}

En effet, $2)\Rightarrow1)$ en vertu de A), et pour $1)\Rightarrow2)$ il suffit de prendre $K=Q(a,b)$ où $a=x+\bar x$, $b=x\bar x$, et appliquer encore A). \uline{Remarque} Cet \unsure{ensemble par les} conditions de 1) est stable par conjugaison, ce qui n'était pas évident autrement.

\deleted{3) $x\bar x$ est tot. str. positif et de même pour $x_\alpha\bar x_\alpha$ pour $\forall$ conjugué de $x$}

C) Soit $x\in\bar Q\subset C$. Conditions équivalentes
\diagram{\begin{tikzcd}[row sep=large]
\begin{gathered}1)\quad x+\bar x\text{ tot. réel,}\quad(x-\bar x)^2\text{ tot. négatif}\end{gathered}\arrow[d,Leftrightarrow]\\
2)\quad\exists K,L\text{ comme plus haut,}\quad x\in\deleted{\unsure{K}}L\arrow[d,Leftrightarrow]\\
3)\quad x\in Q'(i).\arrow[d,Rightarrow,bend right=45]\\
\begin{gathered}\unsure{4)}\quad x\bar x\text{ tot. positif,}\quad\deleted{\unsure{\text{et de même}}}\ \inserted{\text{ainsi que}}\\x_\alpha\bar x_\alpha\text{ pour $\forall$ conjugué $x_\alpha$ de $x$.}\end{gathered}
\end{tikzcd}}
\deleted{\unsure{après}}
\margin{\unsure{2}\ $\Uparrow$}

\mspage{107.}
